\documentclass[11pt,a4paper]{article}

\usepackage[margin=2.2cm]{geometry}
\usepackage{kotex}
\usepackage{amsmath,amssymb}
\usepackage{booktabs}
\usepackage{array}
\usepackage{xcolor}
\usepackage{listings}
\usepackage[most]{tcolorbox}
\usepackage{enumitem}
\usepackage{hyperref}
\usepackage{fancyhdr}
\usepackage{titlesec}

\definecolor{codebg}{RGB}{247,247,247}
\definecolor{codeframe}{RGB}{210,210,210}
\definecolor{keywordcolor}{RGB}{0,70,160}
\definecolor{commentcolor}{RGB}{0,120,80}
\definecolor{stringcolor}{RGB}{170,50,50}
\definecolor{titleblue}{RGB}{35,75,130}

\hypersetup{colorlinks=true,linkcolor=titleblue,urlcolor=titleblue}

\pagestyle{fancy}
\fancyhf{}
\lhead{C++ Programming}
\rhead{Day 6}
\cfoot{\thepage}

\titleformat{\section}{\Large\bfseries\color{titleblue}}{\thesection.}{0.6em}{}
\titleformat{\subsection}{\large\bfseries}{\thesubsection.}{0.5em}{}

\lstdefinestyle{cppstyle}{
    language=C++,
    backgroundcolor=\color{codebg},
    frame=single,
    rulecolor=\color{codeframe},
    basicstyle=\ttfamily\small,
    keywordstyle=\color{keywordcolor}\bfseries,
    commentstyle=\color{commentcolor},
    stringstyle=\color{stringcolor},
    numbers=left,
    numberstyle=\tiny\color{gray},
    numbersep=8pt,
    tabsize=4,
    showstringspaces=false,
    breaklines=true,
    columns=fullflexible
}
\lstset{style=cppstyle}

\newtcolorbox{learningbox}{colback=blue!4,colframe=titleblue,title=\textbf{학습 목표},fonttitle=\bfseries}
\newtcolorbox{importantbox}{colback=yellow!8,colframe=orange!70!black,title=\textbf{핵심 포인트},fonttitle=\bfseries}
\newtcolorbox{practicebox}{colback=red!3,colframe=red!55!black,title=\textbf{실습},fonttitle=\bfseries}

\begin{document}

\begin{titlepage}
    \centering
    \vspace*{3cm}
    {\Huge\bfseries C++ Programming\par}
    \vspace{0.8cm}
    {\LARGE Day 6: 클래스와 객체의 생명주기\par}
    \vspace{2cm}
    {\large 클래스, 멤버 함수, 생성자, \texttt{this}, 정적 멤버, 소멸자\par}
    \vfill
\end{titlepage}

\tableofcontents
\newpage

\section{학습 목표}
\begin{learningbox}
이 장을 마치면 다음 내용을 설명하고 직접 코드로 작성할 수 있다.
\begin{itemize}[leftmargin=1.5em]
    \item 클래스를 정의하고 객체를 생성하기
    \item 멤버 변수와 멤버 함수 사용하기
    \item 생성자로 객체를 초기화하기
    \item \texttt{this} 포인터로 현재 객체의 멤버에 접근하기
    \item 정적 멤버를 사용하여 모든 객체가 값을 공유하게 하기
    \item 소멸자의 호출 시점과 객체의 생명주기 이해하기
\end{itemize}
\end{learningbox}

\section{클래스와 객체}
클래스는 관련된 데이터와 기능을 하나로 묶어 새로운 자료형을 만드는 기능이다. 클래스는 설계도이고, 그 설계도로 실제로 만든 변수를 객체라고 한다.

\begin{lstlisting}
#include <iostream>
#include <string>

using namespace std;

class Student
{
public:
    string name;
    int age;
};

int main()
{
    Student student;

    student.name = "Alice";
    student.age = 20;

    cout << student.name << endl;
    cout << student.age << endl;

    return 0;
}
\end{lstlisting}

\begin{importantbox}
\texttt{public} 아래의 멤버는 클래스 밖에서도 점 연산자 \texttt{.}를 사용해 접근할 수 있다. 클래스는 기본 접근 지정자가 \texttt{private}이므로, 외부 접근이 필요한 멤버에는 \texttt{public}을 명시해야 한다.
\end{importantbox}

\section{멤버 함수}
클래스 안에 정의된 함수를 멤버 함수라고 한다. 멤버 함수는 같은 객체의 멤버 변수에 직접 접근할 수 있다.

\begin{lstlisting}
#include <iostream>
#include <string>

using namespace std;

class Student
{
public:
    string name;
    int age;

    void introduce()
    {
        cout << "Name: " << name << endl;
        cout << "Age: " << age << endl;
    }
};

int main()
{
    Student student;
    student.name = "Alice";
    student.age = 20;
    student.introduce();

    return 0;
}
\end{lstlisting}

\section{생성자}
생성자는 객체가 생성될 때 자동으로 호출되는 특별한 멤버 함수이다. 클래스 이름과 같은 이름을 사용하고 반환형을 작성하지 않는다.

\begin{lstlisting}
#include <iostream>
#include <string>

using namespace std;

class Book
{
public:
    string title;
    int price;

    Book(string t, int p)
    {
        title = t;
        price = p;
    }

    void showInfo()
    {
        cout << "Title: " << title << endl;
        cout << "Price: " << price << endl;
    }
};

int main()
{
    Book book("C++ Basics", 25000);
    book.showInfo();

    return 0;
}
\end{lstlisting}

\begin{importantbox}
멤버 변수는 객체에 저장되어 객체가 살아 있는 동안 유지된다. 반면 생성자의 매개변수는 생성자가 실행되는 동안만 존재한다.
\end{importantbox}

\section{\texttt{this} 포인터}
\texttt{this}는 현재 멤버 함수를 실행하고 있는 객체 자신의 주소를 가리키는 포인터이다. 매개변수와 멤버 변수의 이름이 같을 때 두 변수를 구분하는 데 사용할 수 있다.

\begin{lstlisting}
#include <iostream>
#include <string>

using namespace std;

class Car
{
public:
    string brand;
    int year;

    Car(string brand, int year)
    {
        this->brand = brand;
        this->year = year;
    }

    void showInfo()
    {
        cout << "Brand: " << this->brand << endl;
        cout << "Year: " << this->year << endl;
    }
};

int main()
{
    Car car("Hyundai", 2025);
    car.showInfo();

    return 0;
}
\end{lstlisting}

\begin{importantbox}
\texttt{this->brand}는 현재 객체의 멤버 변수이고, 오른쪽의 \texttt{brand}는 생성자의 매개변수이다. \texttt{brand = brand;}라고 작성하면 가까운 범위의 매개변수끼리 대입하게 된다.
\end{importantbox}

\section{정적 멤버}
일반 멤버 변수는 객체마다 따로 존재하지만, 정적 멤버는 클래스 전체에서 하나만 존재하며 모든 객체가 공유한다.

\begin{lstlisting}
#include <iostream>
#include <string>

using namespace std;

class Student
{
public:
    string name;
    static int studentCount;

    Student(string name)
    {
        this->name = name;
        studentCount++;
    }

    void showInfo()
    {
        cout << "Name: " << name << endl;
    }
};

int Student::studentCount = 0;

int main()
{
    Student s1("Alice");
    Student s2("Bob");
    Student s3("Charlie");

    s1.showInfo();
    s2.showInfo();
    s3.showInfo();

    cout << "Total Students: "
         << Student::studentCount << endl;

    return 0;
}
\end{lstlisting}

\begin{importantbox}
정적 멤버 변수는 클래스 안에서 선언하고 클래스 밖에서 \texttt{자료형 클래스이름::멤버이름} 형식으로 정의한다. 특정 객체가 아니라 클래스에 속하므로 \texttt{Student::studentCount}와 같이 접근하는 것이 명확하다.
\end{importantbox}

\section{소멸자}
소멸자는 객체가 사라질 때 자동으로 호출되는 특별한 멤버 함수이다. 클래스 이름 앞에 \texttt{\textasciitilde}를 붙이며 반환형과 매개변수가 없다.

\begin{lstlisting}
#include <iostream>
#include <string>

using namespace std;

class Car
{
public:
    string brand;

    Car(string brand)
    {
        this->brand = brand;
        cout << brand << " is created." << endl;
    }

    ~Car()
    {
        cout << brand << " is destroyed." << endl;
    }
};

int main()
{
    Car car("Hyundai");

    return 0;
}
\end{lstlisting}

지역 객체는 자신이 선언된 블록을 벗어날 때 소멸한다. 같은 블록에서 여러 객체가 생성되면 소멸자는 생성의 반대 순서로 호출된다.

\begin{lstlisting}
int main()
{
    Car c1("Hyundai");
    Car c2("Kia");
    Car c3("Tesla");

    return 0;
}
\end{lstlisting}

생성 순서는 \texttt{Hyundai}, \texttt{Kia}, \texttt{Tesla}이고, 소멸 순서는 \texttt{Tesla}, \texttt{Kia}, \texttt{Hyundai}이다.

\begin{importantbox}
스택에 생성된 지역 객체는 블록을 벗어나면 자동으로 소멸한다. \texttt{new}로 힙에 생성한 객체는 \texttt{delete}를 호출해야 소멸자가 실행되고 메모리가 해제된다.
\end{importantbox}

\section{객체의 생명주기}
객체는 생성되고, 사용된 뒤, 소멸한다.

\begin{enumerate}[leftmargin=2em]
    \item 객체를 생성하면 생성자가 자동으로 호출된다.
    \item 객체의 멤버 변수와 멤버 함수를 사용한다.
    \item 객체가 유효 범위를 벗어나면 소멸자가 자동으로 호출된다.
    \item 소멸자 실행이 끝난 뒤 객체가 차지하던 메모리가 회수된다.
\end{enumerate}

생성자는 초기 설정을 담당하고, 소멸자는 객체가 사라지기 전 마지막 정리를 담당한다.


\section{실습 목록}
\begin{practicebox}
Day 6 실습 파일은 다음과 같이 관리한다.
\begin{itemize}[leftmargin=1.5em]
    \item \texttt{practice/p01\_class.cpp}
    \item \texttt{practice/p02\_member\_function.cpp}
    \item \texttt{practice/p03\_constructor.cpp}
    \item \texttt{practice/p04\_this\_pointer.cpp}
    \item \texttt{practice/p05\_static\_member.cpp}
    \item \texttt{practice/p06\_destructor.cpp}
\end{itemize}
\end{practicebox}


\section{핵심 요약}
\begin{importantbox}
\begin{itemize}[leftmargin=1.5em]
    \item 클래스는 데이터와 기능을 하나로 묶은 사용자 정의 자료형이다.
    \item 객체는 클래스를 바탕으로 실제로 생성한 변수이다.
    \item 멤버 함수는 같은 객체의 멤버 변수에 직접 접근할 수 있다.
    \item 생성자는 객체 생성 시 자동 호출되며 멤버를 초기화한다.
    \item \texttt{this}는 현재 객체의 주소를 가리키는 포인터이다.
    \item 정적 멤버는 클래스 전체에서 하나만 존재하며 모든 객체가 공유한다.
    \item 소멸자는 객체가 사라질 때 자동으로 호출된다.
    \item 같은 범위의 지역 객체는 생성의 반대 순서로 소멸한다.
\end{itemize}
\end{importantbox}

\end{document}

